%=====================================================
\section{Relations}
%=====================================================

%-----------------------------------------------------
\begin{frame}{Binary Relation}

\textbf{Definition}

A binary relation from $A$ to $B$ is any subset of

\[
A\times B
\]

Mathematically,
\[
R\subseteq A\times B \;if \;(a,b)\in R\]
We write aRb meaning "$a$ is related to $b$."

\end{frame}

%-----------------------------------------------------

\begin{frame}{Example of a Binary Relation}

Let
\[
A=\{1,2,3\}
\]

Define
\[
R=
\{
(1,2),
(2,3),
(1,3)
\}
\]

Then
\[
R\subseteq A\times A
\]

Hence,
\[
R\;is\; a\; binary\; relation\; on\; A.
\]

\end{frame}

%-----------------------------------------------------

\begin{frame}{Predicate Representation}

A relation is often defined using a predicate.

\vspace{0.4cm}

Example

\[
R=
\{
(a,b)
\mid
a,b\in A
\land
a<b
\}
\]

where

\[
P(a,b):a<b
\]

is the predicate.

\end{frame}

%-----------------------------------------------------

\begin{frame}{Domain and Range of a Relation}

For

\[
R\subseteq A\times B
\]

the \textbf{Domain} is

\[
Dom(R)
=
\{
a\in A
\mid
\exists b\in B,
(a,b)\in R
\}
\]

The \textbf{Range} is

\[
Ran(R)
=
\{
b\in B
\mid
\exists a\in A,
(a,b)\in R
\}
\]

\end{frame}